\scanpage{7}
\uline{Proposition} 1 \quad Soit $A$ un anneau noeth. $\unsure{I|J}$-adique séparé complet, $M$\deleted{\unsure{odules|$N$}}$,N$ deux $A$-modules de type fini, et $\mathfrak X=\operatorname{Spf} A$, $M,N$ les Modules coh sur $\mathfrak X$ définis par $M,N$. Alors on a des isomorphismes
\[
 \operatorname{Ext}^{\bullet}(\unsure{x|X|\mathfrak X};M,N)\simeq\operatorname{Ext}^{\bullet}_A(M,N)
\]

\uline{Corollaire} 2 \quad On a
\[
 \underline{\operatorname{Ext}}^{\bullet}(M,N)\simeq\operatorname{Ext}^{\bullet}_A(M,N)^{\triangle}\ \deleted{\unsure{\text{formel}}}
\]

[en effet, on sait que \inserted{les} $\underline{\operatorname{Ext}}^{i}(M,N)$ est cohérent, donc par la suite spectrale et la nullité de la cohomologie des cohérents sur $\mathfrak X$, on a $\operatorname{Ext}^{\bullet}_{\deleted{\unsure{A|\mathfrak X}}}(\mathfrak X;M,N)\simeq H^0(\mathfrak X,\underline{\operatorname{Ext}}^{\bullet}(M,N))$ \deleted{et comme $\operatorname{Ext}^{\bullet}(\mathfrak X;M,N)\simeq\operatorname{Ext}^{\bullet}_A(M,N)$, on conclut} d'où $H^0(\mathfrak X,\underline{\operatorname{Ext}}^{\bullet}(M,N))\simeq\operatorname{Ext}^{\bullet}_A(M,N)$ par \unsure{prp 1|prop 1}, d'où le cor 2 comme $\underline{\operatorname{Ext}}^{\bullet}(M,N)$ est coh.]

Il suffit en effet de prendre une résolution \deleted{finie} de $M$ par des modules libres de type fini pour trouver le résultat. Autre procédé : comme $\mathfrak X\longrightarrow X=\unsure{\operatorname{Sp}|\operatorname{Spec}} A$ est plat, \deleted{\unsure{le}} le \unsure{foncteur|foncteurs} image inverse commute à la formation des $\underline{\operatorname{Ext}}^{\bullet}$, d'où
\[
\begin{aligned}
 \underline{\operatorname{Ext}}^{\bullet}_{\underline O_{\mathfrak X}}(M,N)
 &\simeq\underline{\operatorname{Ext}}^{\bullet}_{\underline O_X}(\widetilde M,\widetilde N)^{\wedge},\quad
 \text{or }\underline{\operatorname{Ext}}^{\bullet}(\widetilde M,\widetilde N)\\
 &=\operatorname{Ext}^{\bullet}_A(M,N)^{\sim}\quad\unsure{(\text{cf})|(\text{cq})},\quad\text{d'où}\\
 &\underline{\operatorname{Ext}}^{\bullet}_{\underline O_{\mathfrak X}}(M,N)\simeq\operatorname{Ext}^{\bullet}_A(M,N)^\triangle
\end{aligned}
\]
(c'est le cor. 2) et en prenant les sections, on trouve à rebours la prop. 1.
% archive sheet number: 6

